\documentclass[12pt]{article}

\usepackage[a4paper, total={6in, 8in}]{geometry}
\usepackage{textcomp}

\begin{document}
\setlength{\parindent}{0pt}

\def \t {\textrightarrow}
\def \v {\vspace{1em}}
\def \bi {\begin{itemize}}
\def \ei {\end{itemize}}
\def \s[#1] {\section*{#1}}
\def \ss[#1] {\subsection*{#1}}
\def \sss[#1] {\subsubsection*{#1}}

\s[Novecento]
Novecento è un gruppo di artisti e architetti che si raccoglie attorno alla figura di Margherita Sarfatti \t a Milano, lei è una critica d'arte \\
Viene da famiglia venenziana ebrea? \t sara per molti anni amante di mussolini, che conosce alla redazione dell'Avanti \\
Relazione finira, e lei andra via dall'italia con le leggi razziali

\v

Questi artisti si pongono l'obbiettivo di ricreare un arte italiana \t in realta c'è fenomeno degli anni 20 chiamato "Ritorno all'ordine" \\
Dopo 1 guerra mondiale europa è distrutta psicologicamente \t c'è bisogno di certezze e serenità, che non si trova nelle avanguardie, che esprimono la crisi delle certezze di primo novecento \t non rasserenano \\
Si vede quindi ritorno al naturalismo, alla prospettiva etc.

\ss[L'allieva - Sironi]
Sironi è futurista, è uno dei protagonisti \t quest'opera è del 24 \t ritratto con figura femminile, statica \\
Forte chiaro/scuro \t e dietro c'è una statuetta che è un nudo che rimanda alla classicita \\
Sironi sara anche principale esponente dell'arte figurativa fascista

\ss[Donne che corrono sulla spiaggia - Picasso]
Lui ha percorso artistico molto lungo che comprende molte fasi \t reinterpreta novita, etc. \\
In questo periodo ha periodo neoclassico \t anche lui partecipa in questo ritorno all'ordine \\
La destrutturazione cubista non è presente \t sono due figure femminili (sproporzionate e deformate, ma lo stesso che rimandano a un naturalismo)

\v

Si vuole creare quindi un arte fascita e italiana, e ci si rifa quindi al rinascimento e alla classicita \t ma non vuol dire imitazione, solo riferimento \\
Dal punto di vista architettonico \t stesso discorso: si riferisce alla romanita e alla classicita, ma non c'è imitazione (no neoclassicismo, ma un utilizzo degli elementi del linguaggio classico che vengono pero semplificati) \\
Abbiamo Portalupi, Giò Ponti, etc. \t fondata rivista "Domus", diretta da Giò Ponti

\sss[Palazzo dell'Arte - Giovanni Muzio]
Ha una funzione espostiva, e viene realizzato da Muzio quando l'esposizione triennale nel 33 viene trasferita a Milano (prima Monza) \\
Palazzo espositivo, che utilizza il linguaggio Novecento \t parte posteriore geometrica ed essenziale, con porticato a forma di C verso il giardino \\
È una successione di arcata a tutto sesto \t che rimanda alla romanita, ma è estremamente essenziale \t non ci sono capitelli, colonne etc.

\sss[Planetario - Piero Portalupi]
È una struttura che viene costruita da Piero Portalupi e viene donata alla città da Hoepli \\
Viene ripresa la classicità puntualmente \t la struttura rimanda al Pantheon, c'è pronao di accesso con i caratteri neoclassici + parallelepipedo di raccordo, e la sala con la cupola (in calcestruzzo però)

\sss[La casa della meridiana]
In via Marchioni (vicino al Pini) \t è un condominio alto-borghese particolare: viene costruito con attenzione alla funzionalità \\
Ha 6 piani, che sono come delle ville sovrapposte \t ogni piano ha appartamento con terrazza, e le terrazze sono poste in modo da garantire la privacy \\
L'ascensore entra direttamente negli appartamenti \\
Funzionalità e tecnologia lo avvicinano al razionalismo, ma marmo + cornici aggettanti (decorative quasi) + grande meridiana ricollegano al novecento

\ss[Ca' Bruta - Giovanni Muzio]
Si trova all'angolo tra via Turati e via Moscova \t è un isolato grande, con forma irregolare \\
Si vede una scelta della planimetria che rimanda al razionalismo: non si limita a costruire gli edifici sul perimetro con corte interna (tipica), ma costruisce due corpi di fabbrica:
\bi
  \item uno curvo??
  \item uno lineare
\ei
Cosi si crea anche strada interna privata \t con accesso che rimanda all'arco di trionfo:
\bi
  \item ha 3 fornici, uno solo centrale
  \item rimanda alla Seriana (rinascimentale) per ??
  \item rimanda a linguaggio classico ma senza decorazioni \t è una geometrizzazione semplificata di un arco del trionfo
  \item rivestito in travertino, molto resistente e usato a Roma
\ei
L'edificio è un edificio alto-borghese, e la struttura è in calcestruzzo armato (che ormai si usa in qualsiasi contesto) con copertura in fibro-cemento \\
Fibro-cemento è materiale nuovo e resistente, ma ha l'amianto \t viene usato molto nelle strutture industriali per le coperture e per le tubature fino agli anni 60 

\v

Estremamente moderno:
\bi
  \item riscaldamento centralizzato
  \item acqua corrente
  \item 10 ascensori
  \item garage sotterraneo con i montacarichi
\ei

\sss[Trattamento dei prospetti]
Non c'è facciata principale \\
Ci sono 3 fasce:
\bi
  \item fascia inferiore comprende di solito 2 piani, qua 3 \t è senza decorazioni ed è rivestita in travertino
  \item fascia centrale \t comprende di solito 3, qua solo 2 \t intonacata di grigio scuro \t e ci sono decorazioni liberamente distribuite: linguaggio classico geometrizzato e ripreso nelle forme: schema dell'arco di trionfo simile a quello dell'entrata, tondo, \textit{opus reticulatus} romano
  \item ultima fascia \t decorazioni e volumetrie articolate \t ci sono delle sporgenze e rientranze \t sempre con stessa logica di ripresa della classicità
\ei
I milanesi di allora percepiscono le decorazioni come poste a caso (sono liberamente distribuite, per questo la chiamano Ca' Bruta) \t decorazioni non si rifanno a qualche esempio della classicià, e non hanno un criterio in cui vengono distribuite \\
Per questo non viene compreso il disegno compositivo della facciata \t e viene considerata quindi brutta

\v

Rimandano un po' a certi dipinti metafisici di De Chirico \t nel 1918 con Carrà e il fratello fondano a Ferrara il movimento Metafisica \t a Ferrara erano soldati \\
È un movimento di avanguardia fino a un certo punto \t come prima impressione è anche naturalistico, ed inserisce richiami classici \t però il tutto viene composto in modo da creare un senso di ambiguità \\
Non si comprende per esempio la spazialità \t e De Chirico chiama questo l'"enigma" \t il dipinto deve porre delle domande

\sss[Piazza d'italia - De Chirico]
Due edifici laterali, una statua classiceggiante in mezzo, due persone, un muretto, un treno e un edificio frontale, e dietro delle montagne (dal punto di vista descrittivo) \\
Gli edifici, con successione di arcate \t rimandano al portico della triennale per esempio, ma sono essenziali senza decorazione \\
L'edificio rimanda alla struttura a \textit{tolos} con presistasi e cella \\
Dal punto di vista formale, lo spazio è sì prospettico, dove però la prospettiva è ambigua \t non è costruita geometricamente, e le proporzioni tra gli elementi sono spiazzanti \\
I due uomoini danno ordine di grandezza \t la statua, se si confrotano le proporzioni, è enorme \t e anche l'edificio in fondo, rispetto a loro, è enorme \\
C'è un'ambiguità dimensionale \\
Le ombre sono molto accentuate, e poi è presente un treno \t che è particolare, e si lega alla figura di De Chirico \t nel dipinto lui inserisce anche se stesso \\
Lui ha passione per la classicità (nasce infatti in Grecia, dove passa infanzia quindi ama grecità) \t il padre è un ingegnere, e costruisce la prima rete ferroviaria greca \t per questo il treno
\end{document}
